\section{Numerical results}
\label{sec:numerical}
This section presents numerical experiments designed to validate the theoretical analysis and evaluate the practical performance of the proposed algorithm. 
Specifically, we verify its core properties: convergence, bound-preserving property, and energy stability in \Cref{subsec:test}. 
Numerical results in \Cref{subsec:comparison} benchmark its efficiency and robustness by comparing \Cref{alg:ADMM} with a standard solver \Cref{alg:Newton}. 
In \Cref{subsec:3d}, we test its effectiveness and robustness in three-dimensional settings.

\subsection{Convergence test}
\label{subsec:test}
In this numerical experiment, the parameters are set as follows:
\begin{equation}
    L=2.0, N=256,\epsilon=0.05, \tau=0.01, T=1.0, \theta=4.0.
\end{equation}
The initial condition is
\begin{equation}
    u_0(x,y)=\frac{1+\sin(2\pi x)\sin(2\pi y)}{3}.
\end{equation}
\par For the iterative solver, the penalty parameter $\rho$ is dynamically updated based on the primal and dual residuals to ensure balanced convergence. \\
The convergence criterion is set as
\begin{equation}
    \max\{r,s\}\leq \gamma=10^{-8},
\end{equation}
where $r$ and $s$ denote the primal residual and the dual residual, respectively.
\par All numerical results obtained from the experiment are presented in graphical form in \Cref{fig:1}.
From \Cref{fig:11}, we see the number of ADMM iterations at each time step, which speaks to the algorithm's effectiveness. 
We observe that, except for the initial steps, the iteration count is essentially uniform across all time steps.
\Cref{fig:12} traces the maximum and minimum of $u$ over time, both staying safely within the physical bounds $(0,1)$, demonstrating the bound-preserving property. 
The discrete energy, plotted in \Cref{fig:13}, decays monotonically as expected, confirming energy stability. 
\Cref{fig:14} displays the primal residual during the first time step; 
its linear decay on a semi-log scale indicates a linear convergence rate.
Finally, \Cref{fig:15} captures the dynamial evolution process of the phase variable at different time steps.
\begin{figure}[htbp]
    \centering
    \begin{subfigure}[b]{0.49\textwidth}
        \centering
        \includegraphics[width=\textwidth]{figure/admm_iterations.png}
        \caption{The number of iterations}
        \label{fig:11}
    \end{subfigure}
    \hfill
    \begin{subfigure}[b]{0.49\textwidth}
        \centering
        \includegraphics[width=\textwidth]{figure/bounds_preservation.png}
        \caption{Bound preservation}
        \label{fig:12}
    \end{subfigure}
    
    \vspace{0.5cm}
    
    \begin{subfigure}[b]{0.49\textwidth}
        \centering
        \includegraphics[width=\textwidth]{figure/energy_curve.png}
        \caption{Energy stability}
        \label{fig:13}
    \end{subfigure}
    \hfill
    \begin{subfigure}[b]{0.49\textwidth}
        \centering
        \includegraphics[width=\textwidth]{figure/admm_convergence.png}
        \caption{Primal residual}
        \label{fig:14}
    \end{subfigure}
    \vspace{0.5cm}

    \begin{subfigure}[b]{1.00\textwidth}
        \centering
        \includegraphics[width=\textwidth]{figure/admm_evolution_test_1.png}
        \caption{Dynamical evolution}
        \label{fig:15}
    \end{subfigure}
    
    \caption{Convergence test}
    \label{fig:1}
\end{figure}

\subsection{Performance comparison with Newton's method}
\label{subsec:comparison}
In this numerical experiment, we compare the performance of \Cref{alg:ADMM} with that of \Cref{alg:Newton}. 
\par The following parameters are adopted in this example:
\begin{equation}
L=2.0, N=1024,\epsilon=0.05, \tau=1.0, T=50.
\end{equation}
The parameter $\theta$ is chosen as the variable parameter in this experiment, which is taken as $\theta=3.0,4.0,5.0$. 
Notice that the time step $\tau=1.0$ is significantly large, 
which verifies that our method has no requirement for the time step. \\
The initial condition is
\begin{equation}
    u_0(x,y)=0.01+0.98\text{rand}(x,y).
\end{equation}
where $\text{rand}(x, y)$ is the function producing the random numbers in $(0,1)$. 
\par We set the convergence criterion for \Cref{alg:ADMM} as:
\begin{equation}
\max\{r,s\}\leq \gamma=10^{-6}.
\end{equation}
Correspondingly, the convergence criterion for \Cref{alg:Newton} is
\begin{equation}
\|u^{(k+1)}-u^{(k)}\|_h \leq \gamma=10^{-6}.
\end{equation}
We present the numerical results with \Cref{ta:comparison} and \Cref{fig:2}.
\begin{table}
    \centering
    \begin{tabular}{c|c|c|c|c}
    Algorithm&$\theta$&CPU\_time(s)&Final\_energy&Final\_residual\\
    \hline
    1&3.0&382&-0.109&3.7E-5\\
    \hline
    2&3.0&478&-0.109&6.4E-5\\  
    \hline
    1&4.0&371&0.198&1.3E-4\\
    \hline
    2&4.0&604&0.198&2.0E-4\\
    \hline
    1&5.0&794&0.385&5.5E-4\\
    \hline
    2&5.0&>2000&-&-\\

    \end{tabular}
    \caption{Comparison of algorithm performance}
    \label{ta:comparison}
\end{table}

As indicated in the table, \Cref{alg:ADMM} achieves higher efficiency than \Cref{alg:Newton}. 
Furthermore, it is more robust with respect to the parameter $\theta$ compared to its counterpart.
We also observe that the solution time increases significantly for large values of $\theta$, since the problem becomes increasingly stiff.
The dynamic evolution of the phase variable is depicted in \Cref{fig:2} for different choices of $\theta$.

\begin{figure}[htbp]
    \centering
    \begin{subfigure}[b]{0.78\textwidth}
        \centering
        \includegraphics[width=\textwidth]{figure/admm_evolution_test_23.png}
        \caption{$\theta=3.0$}
        \label{fig:23}
    \end{subfigure}
    \hfill
    
    \vspace{0.5cm}
    
    \begin{subfigure}[b]{0.78\textwidth}
        \centering
        \includegraphics[width=\textwidth]{figure/admm_evolution_test_24.png}
        \caption{$\theta=4.0$}
        \label{fig:24}
    \end{subfigure}
    \hfill

    \vspace{0.5cm}

    \begin{subfigure}[b]{0.78\textwidth}
        \centering
        \includegraphics[width=\textwidth]{figure/admm_evolution_test_25.png}
        \caption{$\theta=5.0$}
        \label{fig:25}
    \end{subfigure}
    
    \caption{Dynamical evolution}
    \label{fig:2}
\end{figure}


\subsection{Three-dimensional simulation}
\label{subsec:3d}
We present a numerical simulation experiment in a three-dimensional setting here. The parameters we adopted are 
\begin{equation}
L=1.0, N=64,\theta=4.0, \tau=0.01, T=0.5.
\end{equation}
The parameter $\epsilon$ is chosen as the variable parameter in this experiment, which is taken as $\epsilon=0.05,0.10,0.15$. \\
The initial condition is
\begin{equation}
    u_0(x,y,z)=0.01+0.98\text{rand}(x,y,z).
\end{equation}
where $\text{rand}(x,y,z)$ is the function producing the random numbers in $(0,1)$. 
\par The convergence criterion is:
\begin{equation}
\max\{r,s\}\leq \gamma=10^{-6}.
\end{equation}

\begin{figure}[htbp]
    \centering
    \begin{subfigure}[b]{0.78\textwidth}
        \centering
        \includegraphics[width=\textwidth]{figure/isosurface_ep=0.05.png}
        \caption{$\epsilon=0.05$}
        \label{fig:31}
    \end{subfigure}
    \hfill
    
    \vspace{0.5cm}
    
    \begin{subfigure}[b]{0.78\textwidth}
        \centering
        \includegraphics[width=\textwidth]{figure/isosurface_ep=0.10.png}
        \caption{$\epsilon=0.10$}
        \label{fig:32}
    \end{subfigure}
    \hfill

    \vspace{0.5cm}

    \begin{subfigure}[b]{0.78\textwidth}
        \centering
        \includegraphics[width=\textwidth]{figure/isosurface_ep=0.15.png}
        \caption{$\epsilon=0.15$}
        \label{fig:33}
    \end{subfigure}
    
    \caption{Dynamical evolution of isosurface (isovalue $u=0.50$)}
    \label{fig:3}
\end{figure}
A fixed isovalue of $u=0.50$ is used to identify the interface between the two phases. 
Shown in \Cref{fig:3} are the isosurface evolution at different time steps from the three-dimensional simulation. 
The results indicate that the algorithm accurately captures the dynamical process of the two-phase interface with satisfactory precision,
which is consistent with the physical phenomenon of phase separation. In summary, the numerical result demonstrates the robustness of our solver in numerical simulation.

\subsection{Adaptive time-stepping simulation}\label{subsec:Adaptive}
In this experiment, we adopt an adaptive time-stepping strategy to demonstrate the high efficiency of the algorithm. 
The parameters we adopted are set as
\begin{equation}
L=2.0, N=256,\epsilon = 0.05, \tau^0=1.0,\theta=4.0, T=10^{12}.
\end{equation}
\par The initial condition is
\begin{align}
u_0(x,y)=
\begin{cases}
    10^{-5},&\text{if}\quad |x-1|\leq 0.35, |y-1|\leq 0.35,\\
    1-10^{-5},&\text{otherwise}
\end{cases}
\end{align}
which is taken from \cite{wang2020stabilized}.
Numerical experiments indicate that the two-phase evolution under this initial condition requires long-time simulation. 
Hence, we employ this example to perform adaptive time-stepping simulation, 
which dynamically adjusts the time step based on the energy decrement to maintain accuracy while reducing computational cost.
To be more precise, the adaptive adjustment procedure is:
\begin{align}
\tau^{n+1}=
\begin{cases}
\min \{\beta_1 \tau^{n},\tau_{\max}\},&\text{if} \quad |\frac{E_h^{n-1}-E_h^n}{E_h^{n-1}}|<\gamma_1,\\
\max \{\beta_2 \tau^{n},\tau_{\min}\},&\text{if} \quad |\frac{E_h^{n-1}-E_h^n}{E_h^{n-1}}|>\gamma_2,\\
\tau^n,&\text{otherwise},
\end{cases}
\end{align}
where $\beta_1=2.0,\tau_{\max}=10^{10},\gamma_1=0.1,\beta_2=0.5,\tau_{\min}=10^{-3},\gamma_2=0.2$.
\par The numerical results are shown in \Cref{fig:4}. 
The same observations as in \Cref{subsec:test} hold here, that is, 
the results confirm linear convergence, energy stability, and the bound-preserving property of the scheme from \Cref{fig:41} to \Cref{fig:44}.
The evolution of the phase field variable over time is shown in \Cref{fig:45}, where the square-shaped phase is seen to rapidly shrink into a circular one.

\begin{figure}[htbp]
    \centering
    \begin{subfigure}[b]{0.49\textwidth}
        \centering
        \includegraphics[width=\textwidth]{figure/admm_iterations_adaptive.png}
        \caption{The number of iterations}
        \label{fig:41}
    \end{subfigure}
    \hfill
    \begin{subfigure}[b]{0.49\textwidth}
        \centering
        \includegraphics[width=\textwidth]{figure/bounds_preservation_adaptive.png}
        \caption{Bound preservation}
        \label{fig:42}
    \end{subfigure}
    
    \vspace{0.5cm}
    
    \begin{subfigure}[b]{0.49\textwidth}
        \centering
        \includegraphics[width=\textwidth]{figure/energy_curve_adaptive.png}
        \caption{Energy stability}
        \label{fig:43}
    \end{subfigure}
    \hfill
    \begin{subfigure}[b]{0.49\textwidth}
        \centering
        \includegraphics[width=\textwidth]{figure/admm_convergence_adaptive.png}
        \caption{Primal residual}
        \label{fig:44}
    \end{subfigure}
    \vspace{0.5cm}

    \begin{subfigure}[b]{1.00\textwidth}
        \centering
        \includegraphics[width=\textwidth]{figure/admm_evolution_test_4.png}
        \caption{Dynamical evolution}
        \label{fig:45}
    \end{subfigure}
    
    \caption{Adaptive time-stepping simulation}
    \label{fig:4}
\end{figure}

